\documentclass[aspectratio=169,10pt,spanish]{beamer}

\input{../../../_preambulo_beamer.tex}
\input{../../../_comandos_beamer.tex}

\selectlanguage{spanish}

\title{MA1001B - Modelación Estadística para la Toma de Decisiones}
\subtitle{Lección 18: Tiempos de espera geométricos y binomiales negativos}
\author{}
\institute{Tecnol\'ogico de Monterrey}
\date{}

\begin{document}

\begin{frame}[plain]
  \vspace{1.2cm}
  {\Large\bfseries MA1001B - Modelación Estadística para la Toma de Decisiones\par}
  \vspace{0.7cm}
  {\large Lección 18: Tiempos de espera geométricos y binomiales negativos\par}
  \vspace{0.7cm}
  {\color{TecRojo}\rule{\textwidth}{1.2pt}}
  \vspace{0.8cm}

  Facultad MA1001B\\[0.3cm]
  Periodo\\[0.3cm]
  Tecnol\'ogico de Monterrey
\end{frame}

\begin{frame}{Esperar el primer éxito}
  \begin{alertblock}{Pregunta guía}
    Si cada ensayo independiente de auditoría es defectuoso con chance
    $p=0.12$, ¿cuánto se espera el primer éxito, y qué ocurre si la fatiga
    baja $p$?
  \end{alertblock}

  \vspace{0.6em}
  Al final de esta lección, usted puede:
  \begin{itemize}
    \item calcular $P(X=1)=p$ y $E[X]=1/p$;
    \item extender la espera a $r=3$ éxitos con $E[X]=r/p$;
    \item convertir el conteo de fallos de scipy en un conteo de ensayos;
    \item explicar por qué un $p$ declinante hace que la espera media exceda $1/p$.
  \end{itemize}

  \vfill
  \scriptsize Todos los tiempos de espera usados hoy son sintéticos. No se usan datos reales de empresa.
\end{frame}

\begin{frame}{Ensayos geométricos hasta el primer defectuoso}
  \small
  $X$ es el número de ensayo del primer éxito, $X=1,2,3,\ldots$
  \[
    P(X=1)=p=0.12,\qquad E[X]=\frac{1}{p}=8.333333.
  \]
  Verificado $P(X\le 5)=0.472268$. El comando \texttt{geom} de scipy usa esta
  convención de ensayos hasta el primer éxito.
\end{frame}

\begin{frame}[fragile]{Paso 1: Tiempo de espera geométrico}
  \begin{columns}[T]
    \begin{column}{0.42\textwidth}
      \small
      \begin{lstlisting}
p = 0.12
p1 = stats.geom.pmf(1, p)
mean = 1 / p
      \end{lstlisting}

      \begin{block}{Salida verificada}
        $P(X=1)=0.120000$\\
        $P(X\le 5)=0.472268$\\
        $E[X]=8.333333$
      \end{block}
    \end{column}
    \begin{column}{0.58\textwidth}
      \centering
      \includegraphics[width=\textwidth]{../figuras/tiempos_espera_01_geometrica.png}
      \vspace{0.2em}
      \scriptsize FMP geométrica del Paso 1
    \end{column}
  \end{columns}
\end{frame}

\begin{frame}{Binomial negativa: esperar $r=3$}
  \small
  Ahora $X$ es el ensayo del tercer éxito.
  \[
    E[X]=\frac{r}{p}=\frac{3}{0.12}=25,\qquad P(X=3)=p^3=0.001728.
  \]
  El comando \texttt{nbinom} de scipy cuenta fallos antes de $r$ éxitos. Los
  scripts suman $r$ para que $X$ sea un conteo de ensayos, coincidiendo con
  OpenStax.
\end{frame}

\begin{frame}[fragile]{Paso 2: Ensayos hasta $r=3$}
  \begin{columns}[T]
    \begin{column}{0.40\textwidth}
      \small
      \begin{lstlisting}
r, p = 3, 0.12
fail = x - r
mass = stats.nbinom.pmf(
    fail, n=r, p=p
)
      \end{lstlisting}

      \begin{block}{Salida verificada}
        $E[X]=25.000000$\\
        $\Var(X)=183.333333$\\
        $P(X=3)=0.001728$
      \end{block}
    \end{column}
    \begin{column}{0.60\textwidth}
      \centering
      \includegraphics[width=\textwidth]{../figuras/tiempos_espera_02_binomial_negativa.png}
      \vspace{0.2em}
      \scriptsize FMP binomial negativa del Paso 2
    \end{column}
  \end{columns}
\end{frame}

\begin{frame}{Paso 3: La fatiga cambia $p$}
  \begin{columns}[T]
    \begin{column}{0.42\textwidth}
      \small
      Sea $p_t=0.12\times 0.97^{t-1}$. Entonces $p_1=0.120000$ pero
      $p_{10}=0.091228$.

      \vspace{0.4em}
      Espera media con semilla $42$: $17.162800$, frente a geométrica
      $8.333333$.

      \vspace{0.4em}
      \textbf{Límite:} los ensayos no son independientes si la fatiga cambia
      $p$.
    \end{column}
    \begin{column}{0.58\textwidth}
      \centering
      \includegraphics[width=\textwidth]{../figuras/tiempos_espera_03_dependencia_fatiga.png}
      \vspace{0.2em}
      \scriptsize Espera media fatigada del Paso 3
    \end{column}
  \end{columns}
\end{frame}

\begin{frame}{Verificación de conceptos}
  \begin{enumerate}
    \item ¿Por qué $P(X=1)$ es igual a $p$, y no a $(1-p)p$?
    \item Calcule $E[X]$ para el modelo geométrico con $p=0.12$.
    \item ¿Cuál es $E[X]$ si se espera $r=3$ éxitos?
    \item ¿Por qué un $p$ declinante hace la espera más larga que $1/p$?
  \end{enumerate}

  \vspace{0.6em}
  \begin{block}{Conexión con la Lección 19}
    La sesión puede continuar con variables aleatorias continuas, donde las
    probabilidades son áreas y no masas puntuales.
  \end{block}

  \vfill
  \scriptsize Fuente: OpenStax, \textit{Introductory Business Statistics 2e},
  Sección 4.3. CC BY-NC-SA.
\end{frame}

\appendix



\end{document}
